\documentclass{article}

% if you need to pass options to natbib, use, e.g.:
\PassOptionsToPackage{numbers, compress}{natbib}
% before loading neurips_2026

% The authors should use one of these tracks.
% Before accepting by the NeurIPS conference, select one of the options below.
% 0. "default" for submission
% \usepackage{neurips_2026}
% the "default" option is equal to the "main" option, which is used for the Main Track with double-blind reviewing.
% 1. "main" option is used for the Main Track
%  \usepackage[main]{neurips_2026}
% 2. "position" option is used for the Position Paper Track
%  \usepackage[position]{neurips_2026}
% 3. "eandd" option is used for the Evaluations & Datasets Track
 % \usepackage[eandd]{neurips_2026}
% 4. "creativeai" option is used for the Creative AI Track
%  \usepackage[creativeai]{neurips_2026}
% 5. "sglblindworkshop" option is used for the Workshop with single-blind reviewing
 % \usepackage[sglblindworkshop]{neurips_2026}
% 6. "dblblindworkshop" option is used for the Workshop with double-blind reviewing
%  \usepackage[dblblindworkshop]{neurips_2026}

% After being accepted, the authors should add "final" behind the track to compile a camera-ready version.
% 1. Main Track
%  \usepackage[main, final]{neurips_2026}
% 2. Position Paper Track
%  \usepackage[position, final]{neurips_2026}
% 3. Evaluations & Datasets Track
 % \usepackage[eandd, final]{neurips_2026}
% 4. Creative AI Track
%  \usepackage[creativeai, final]{neurips_2026}
% 5. Workshop with single-blind reviewing
%  \usepackage[sglblindworkshop, final]{neurips_2026}
% 6. Workshop with double-blind reviewing
%  \usepackage[dblblindworkshop, final]{neurips_2026}
% Note. For the workshop paper template, both \title{} and \workshoptitle{} are required, with the former indicating the paper title shown in the title and the latter indicating the workshop title displayed in the footnote.
% For workshops (5., 6.), the authors should add the name of the workshop, "\workshoptitle" command is used to set the workshop title.
% \workshoptitle{WORKSHOP TITLE}

% "preprint" option is used for arXiv or other preprint submissions
 \usepackage[preprint]{neurips_2026}

% to avoid loading the natbib package, add option nonatbib:
%    \usepackage[nonatbib]{neurips_2026}

\usepackage[utf8]{inputenc} % allow utf-8 input
\usepackage[T1]{fontenc}    % use 8-bit T1 fonts
\usepackage{hyperref}       % hyperlinks
\usepackage{url}            % simple URL typesetting
\usepackage{booktabs}       % professional-quality tables
\usepackage{amsfonts}       % blackboard math symbols
\usepackage{amsmath}        % math
\usepackage{nicefrac}       % compact symbols for 1/2, etc.
\usepackage{microtype}      % microtypography
\usepackage{xcolor}         % colors

% Note. For the workshop paper template, both \title{} and \workshoptitle{} are required, with the former indicating the paper title shown in the title and the latter indicating the workshop title displayed in the footnote.
\title{Project: Unitree G1 Soccer Skills}


% The \author macro works with any number of authors. There are two commands
% used to separate the names and addresses of multiple authors: \And and \AND.
%
% Using \And between authors leaves it to LaTeX to determine where to break the
% lines. Using \AND forces a line break at that point. So, if LaTeX puts 3 of 4
% authors names on the first line, and the last on the second line, try using
% \AND instead of \And before the third author name.


\author{CS2810 Teaching Team}


\begin{document}

\maketitle

\section{Overview}

This is the project component of Embodied AI CS2810 (Spring 2026). The task is to develop reinforcement learning skills for humanoid robot soccer using the Unitree G1 (29 dof) platform in \texttt{mjlab} physics simulation. The project centers on two tasks:

\begin{itemize}
  \item \textbf{Shooting}: The G1 robot must approach and kick a stationary ball toward a goal. This is a perception-guided motor skill requiring motion tracking priors and ball trajectory prediction.
  \item \textbf{Goalkeeping}: The G1 robot must intercept an incoming ball launched along a parabolic trajectory. This is a reactive interception task requiring rapid end-effector positioning and ball trajectory estimation.
\end{itemize}

\section{Project Phases \& Grading Rubric}
The project is evaluated out of 100 total points, broken down into three primary components: Phase 1 (60\%), Phase 2 (30\% = 15\% Shooting + 15\% Goalkeeping), and Submission Deliverables (10\%).

\subsection{Phase 1: Basic Shooting and Goalkeeping (60\% of project score)}
Implement single-agent training for individual shooter and goalkeeper tasks in a non-adversarial setting. You may refer to \cite{kong2026learningsoccerskillshumanoid} and \cite{ren2025humanoidgoalkeeper} for guidance on environment design, reward shaping, and training strategies. This phase evaluates the fundamental motor control and perception integration of your policies.

The 60\% allocation is divided equally between the two tasks (30\% each), calculated over 50 evaluation trials:
\begin{itemize}
    \item \textbf{Shooter Policy (30 points total):} Evaluated using \texttt{eval\_naive\_shooter.py}. Success is defined as the ball completely crossing the goal plane inside the frame posts.
    \item \textbf{Goalkeeper Policy (30 points total):} Evaluated using \texttt{eval\_naive\_goalkeeper.py}. Success is defined as preventing the ball from entering the goal (ball does not cross the goal plane inside the frame).
\end{itemize}

Both tasks are scored with a linear formula applied to the success/block rate over 50 evaluation trials:
\[
\text{score} = \max\left(0, \frac{\text{success\_rate} - 0.8}{0.2}\right) \times 30
\]
Success rates below 80\% receive 0 points. At 80\% the score is 0, at 100\% the score is 30, with linear interpolation in between.

\textbf{Reproducibility.} The eval scripts (\texttt{eval\_naive\_shooter.py} and \texttt{eval\_naive\_goalkeeper.py}) support a \texttt{-\/-seed} flag (default: 2810) that seeds Python, NumPy, PyTorch, CUDA, and the environment RNGs, making each run fully reproducible. For your Phase~1 report, you must evaluate your policies with \textbf{three different seeds}: \texttt{42}, \texttt{2810}, and \texttt{202606}. Run 50 trials per seed, compute the mean success/block rate across the three runs, and apply the scoring formula above using this averaged rate. Reporting the mean over three independent seeds provides a more stable and fair metric than a single 50-trial run.

\subsection{Phase 2: Cross-Evaluation Tournament (30\% of project score)}

Phase 2 is an all-to-all cross-evaluation tournament conducted centrally by the teaching team. The 30\% allocation is divided equally between shooting and goalkeeping (15\% each). After completing basic single-agent policies in Phase 1, teams continue to improve their agents---through adversarial training, self-play, or further single-agent optimization---and then compete against every other team.

Because each team will design unique observation spaces, feature extractors, and neural network architectures, your codebases will be heterogeneous. To enable cross-evaluation without requiring teams to share code or checkpoints, the tournament adopts a standardized API interface: each team deploys their trained policy behind a lightweight server, and the competition script \texttt{scripts/compete.py} loads two G1 robots into a single MuJoCo scene and queries each team's server at every simulation step.

\subsubsection{API Protocol}

All teams must implement a policy server exposing the following two endpoints (a reference implementation is provided as \texttt{scripts/api\_server.py}):

\begin{itemize}
    \item \texttt{POST /act} --- Receives raw MuJoCo state for both robots and the ball, and returns an action for the robot the server controls (determined by the \texttt{-\/-task} flag).
    \begin{verbatim}
    Request:  {
      "shooter":   {root_pos, root_quat, root_lin_vel, root_ang_vel,
                    joint_pos, joint_vel, last_action},
      "goalkeeper": {...},
      "ball":      {pos, vel}
    }
    Response: {"action": [[f32, ...]]}\end{verbatim}
    \item \texttt{POST /reset} --- Resets the policy's internal state (LSTM hidden state, history buffer, etc.).
    \begin{verbatim}
    Request:  {}
    Response: {"status": "ok"}\end{verbatim}
\end{itemize}

\texttt{compete.py} sends the same raw state to both servers regardless of their role---each server extracts its own robot's data to compute observations.

\subsubsection{Timeline and Milestones}
\begin{itemize}
    \item \textbf{Weeks 13--16: Agent Development \& Improvement} \\
    Teams first complete their basic single-agent shooting and goalkeeping policies (Phase 1). They then continuously improve their agents' capabilities through adversarial training, self-play, or further single-agent optimization. And this phase will continue until the end of week 16 (6.21).

    \item \textbf{Week 17: Tournament \& Report Submission} \\
    The tournament is conducted centrally on a server provided by the teaching team (we will make the server available later). Each team deploys their policy servers (which must be network-reachable from the tournament server) and submits the access URLs. The teaching team runs \texttt{scripts/compete.py} for every pair of teams (Team A and Team B), evaluating both directions:
    \begin{itemize}
        \item Team A's Shooter vs.\ Team B's Goalkeeper (10 trials)
        \item Team B's Shooter vs.\ Team A's Goalkeeper (10 trials)
    \end{itemize}
    The specific evaluation schedule will be announced before week 17. But notice that both the tournament results and the final report are due by the end of Week 17 (6.26).
\end{itemize}

\subsubsection{Scoring Mechanism}
There are 14 teams in total. Each team's Phase 2 score is computed separately for shooting and goalkeeping (15 points maximum each):

\begin{itemize}
    \item \textbf{Base Score:} Every team starts with 2 points for each task.
    \item \textbf{Head-to-Head Points:} For each opponent team, a direct matchup is played. Taking Team A's Shooter vs.\ Team B's Goalkeeper as an example: 10 trials are run. If Team A's shooter scores $>5$ goals, Team A earns +1 point; otherwise (i.e., the goalkeeper holds the shooter to $\le 5$ goals), Team B earns +1 point. The same rule applies symmetrically to Team A's Goalkeeper vs.\ Team B's Shooter.
\end{itemize}

Each team plays every other team once per task direction, yielding 13 head-to-head opportunities per task. The maximum score per task is therefore $2 + 13 = 15$ points.

\subsection{Submission and Deliverables (10\% of project score)}
Submission consists of a PDF report and the corresponding code repository before the end of Week 17 (6.26) to gradescope. The PDF report must include:
\begin{itemize}
    \item \textbf{Phase 1 Results:} Basic shooting and goalkeeping evaluation results (success/block rates over 50 trials) with analysis.
    \item \textbf{Phase 2 Design \& Approach:} Description of how you improved your agents for the tournament (adversarial training, self-play, or further single-agent optimization). Explain your API server implementation, including how you compute observations from raw state. Since the tournament is conducted centrally by the teaching team, you do not need to report cross-evaluation results yourself.
\end{itemize}


\section{Using the Template}

The template is given at \href{https://github.com/xiaojxkevin/G1-mjlab-soccer}{https://github.com/xiaojxkevin/G1-mjlab-soccer}.

\subsection{Repository Structure}

The template provides:

\begin{itemize}
  \item \textbf{Playable environments}: \texttt{Unitree-G1-Shooter} and \texttt{Unitree-G1-Goalkeeper} registered as mjlab tasks with configurable scene layout, ball physics, and MuJoCo visualization.
  \item \textbf{Evaluation scripts}: \texttt{eval\_naive\_shooter.py} and \texttt{eval\_naive\_goalkeeper.py} that load trained checkpoints, run headless batch trials, collect paper-matching metrics, and record videos. And you are allowed to modify these scripts if you want to use different observation spaces, reward functions etc.
  \item \textbf{Training configs}: Training environment designs under \texttt{config/training/} that specify observation spaces, reward functions, termination conditions, and domain randomization from the two papers. And it should be noticed that these are generated with coding agents and may contain errors. So you should carefully review and debug them before use.
\end{itemize}

\subsection{Play and Visualization}

Use the \texttt{scripts/play.py} script to visualize the environment:

\begin{verbatim}
# Shooter
python scripts/play.py Unitree-G1-Shooter --agent zero --viewer native

# Goalkeeper
python scripts/play.py Unitree-G1-Goalkeeper --agent zero --viewer native
\end{verbatim}

\subsection{Evaluation}

Use the eval scripts to benchmark trained policies. See Section~1.1 for the
required seed and trial settings for Phase~1 scoring.

\begin{verbatim}
# Shooter
python scripts/eval_naive_shooter.py --headless --num-trials 50 \
    --checkpoint <path>

# Goalkeeper
python scripts/eval_naive_goalkeeper.py --headless --num-trials 50 \
    --checkpoint <path>
\end{verbatim}

\subsection{Implementing Training}

The \texttt{train.py} script in \texttt{scripts/} serves as the training entrypoint. It is expected to register their own training tasks, implement the model classes, and run PPO training.

\subsection{Cross-Evaluation Competition}

For Phase 2 cross-evaluation, teams deploy their policies as API servers and use \texttt{scripts/compete.py} to run matchups:

\begin{verbatim}
# Start policy servers (one per robot)
python scripts/api_server.py --checkpoint <path> --port 8000 --task shooter
python scripts/api_server.py --checkpoint <path> --port 8001 --task goalkeeper

# Run a matchup
python scripts/compete.py \
    --shooter-api http://<team_a_ip>:8000 \
    --goalkeeper-api http://<team_b_ip>:8001 \
    --headless --num-trials 10
\end{verbatim}

The script loads two G1 robots into a single MuJoCo scene, reads raw MuJoCo state (joint angles, root poses, ball position/velocity), and sends the same state dict to both teams' API servers. Each team's server computes its own observation tensor and returns an action. This design fully decouples observation spaces across teams. A reference API server with example observation computation is provided as \texttt{scripts/api\_server.py}.

\subsection{Contributing and Reporting Issues}

This template is a work in progress. The training configs, environment wrappers, and evaluation scripts were partially generated by coding agents and may contain bugs or rough edges. \textbf{We strongly encourage all students to}:
\begin{itemize}
    \item \textbf{Report issues} --- if you find a bug, a broken config, or unclear documentation, open a GitHub Issue on the template repository.
    \item \textbf{Submit pull requests} --- fixes, improvements, and additional utilities (e.g., better observation configs, visualization tools, or multi-agent helpers) are welcome. PRs that benefit the whole class will be merged and acknowledged.
\end{itemize}

Treat this repository as a living codebase that improves with your feedback throughout the semester.

{
\bibliographystyle{IEEEtran}
\bibliography{ref}
}

\end{document}
